\chapter{\Peano's particle support}
\label{section:pidt}

\noindent
\Peano's particle handling is a follow-up to the 2015 paper ``Two Particle-in-Grid Realisations on Spacetrees'' (\url{https://arxiv.org/abs/1508.02435}).




\begin{framed}
\paragraph*{A first \Peano\ pidt run (native usage of particle interface)}
\begin{itemize}
  \item Configure \Peano\ with the flags \texttt{--enable-particles}:
  \begin{code}
libtoolize; aclocal; autoconf; autoheader; cp src/config.h.in .; automake --add-missing
./configure --enable-particles
  \end{code}
   These are the two mandatory extensions. You might want to add further options.
  \item Recompile your \Peano\ code. This should give you a couple of
  \texttt{ToolboxParticles} libraries in the \texttt{src/toolbox/particles}
  subdirectory.
  \item Ensure that \texttt{python} subdirectory is in your Python search path
  and
  \item run the Jupyter notebook in \texttt{examples/particles/sph}: 
  \begin{code}
export PYTHONPATH=../../../python/
jupyter-notebook peano-sph.ipynb
  \end{code}
%   \noindent
%   or---if you don't like Jupyter notebooks---type in
%   \begin{code}
% export PYTHONPATH=../../../python/
% ipython3 example-scripts/finitevolumes-with-ExaHyPE2-benchmark.py --h 0.1
% ./peano4
%   \end{code}
\end{itemize}
\end{framed}


\noindent
\Peano's particle code is lightweight.
There are some C++ building blocks, but the majority of its features are
available through the Python API as
\texttt{toolbox.particles}.
You will need

\begin{code}
import peano4.toolbox.particles
\end{code}

\noindent
I assume that you are familiar from hereon with \Peano's basics, i.e.~you know 
how to create a project, 


\section{Introduce a new particle}

With the Python API, the generation of a particle is straightforward:

\begin{code}
import dastgen2


particle = peano4.toolbox.particles.Particle("Particle")

v_attr = peano4.dastgen2.Peano4DoubleArray("v","Dimensions")
mass_attr = dastgen2.attributes.Double("mass")

particle.data.add_attribute(v_attr)
particle.data.add_attribute(mass_attr)

particles = peano4.toolbox.particles.ParticleSet(particle)
\end{code}

\begin{remark}
 Though conceptionally a plain C++ pretty printer, the DaStGen2 Python API
 can create some fancy C++ annotations. You can, for example, specify ranges
 in which integers live, or you can prescribe the precision you wanna have 
 for floating point values. If you have a specialised LLVM compiler, these
 annotations will be used to generate aggressively optimised code.
\end{remark}


\noindent
The \Peano\ particle has some additional attributes as they are discussed in 
Section \ref{section:swift:backbone:data-model} and thereafter. 
For the user, it is usually sufficient to recognise that each particle by 
default has a position \texttt{x} and and a cut-off radius
\texttt{cutOffRadius}.
You can query and modify both values through \texttt{setX}, \texttt{getX},
\texttt{getCutOffRadius} and \texttt{setCutOffRadius}, respectively.


We next add the particles to the project and we notify the vertices that there 
are particles which have to be administered. 
I discuss why and how pidt uses vertices in Section
\ref{section:pidt:backbone:dual-trees} though users might (initially) skip such
background info. 

\begin{code}
project.datamodel.add_global_object(particle)
project.datamodel.add_vertex(particles)
\end{code}

\noindent
Once you have added the particle and the particle container \texttt{particles}
to your \Peano\ datamodel, you have to inform all solver steps that they have to
administer the particles.

\begin{code}
create_grid.use_vertex(particles)
\end{code}

\noindent
Let's assume that our project is spanning the namespace
\texttt{examples::someparticles}.
 In the example above where our particle is
called \texttt{Particle}, we will end up with two classes:
\linebreak
\texttt{examples::someparticles::globaldata::Particle} is the actual particle
class. 
When you open the generated C++ code, you'll recognise that it is a 1:1 mapping
of the Python attributes onto a C++ class with some additional routines used by
\Peano\ internally.
There is a further class
\texttt{examples::someparticles::vertexdata::ParticleSet}.
As the name says, this is the actual data associated with each vertex.
It is a specialisation of \texttt{std::vector} over pointers to
\texttt{examples::someparticles::globaldata::Particle}.

We summarise that each vertex holds a vector to pointers over particles.
The actual particles reside on the heap.
From hereon, you introduce, as with any \Peano\ project additional data fields,
action sets, and so forth, and you write you main to actually control the
program flow.
We assume from hereon that the reader is familiar with the concept of an
program phase (algorithmic step) mapped onto an action set.



\section{Insert particles into domain}

To create a particle, you in principle just allocate it on the heap, and then
you add it to its nearest vertex.
Furthermore, there are some mandatory attributes that aways have to be set. 
Therefore, I recommend that you use the toolbox' initialisation routine to set
up particles correctly.


In principle, it does not matter where you create the particle.
You can do so in \texttt{touchCellFirstTime()} or in a cell event. 
I usually use a cell event and I rely on the helper function
\texttt{insertParticleIntoCell} to take care of the vertex association:


\begin{code}
#include "toolbox/particles/ParticleFactory.h"



...::touchCellFirstTime(
  const peano4::datamanagement::CellMarker& marker,
  peano4::datamanagement::VertexEnumerator<...::vertexdata::ParticleSet> fineGridVerticesParticleSet, 
  ...
  peano4::datamanagement::VertexEnumerator<...::vertexdata::ParticleSet> coarseGridVerticesParticleSet,
  ...
) {
examples::someparticles::globaldata::Particle* myNewParticle = 
  new examples::someparticles::globaldata::Particle();
  
// assume x and cutOffRadius are set to something meaningful
toolbox::particles::init( *myNewParticle, x, cutOffRadius );

// here's your user code's initialisation
myNewParticle->setXXX(...);

toolbox::particles::insertParticleIntoCell(
  marker,
  myNewParticle,
  fineGridVerticesParticleSet
);
}
\end{code}


\noindent
There is no need to write such initialisation code yourself. 
\Peano's particle toolbox provides a couple of pre-defined particle generators
that you can adopt to your needs.
Probably the most popular particle seeds are
\begin{itemize}
  \item \texttt{InsertParticlesAlongCartesianLayoutIntoUnrefinedCells}
  \item \texttt{InsertRandomParticlesIntoUnrefinedCells}
\end{itemize}












\section{Further action sets}

There are a few further action sets in the toolbox that help you to write your
Lagrangian code:


\begin{enumerate}
  \item The action set \texttt{toolbox.particles.PlotParticlesInVTKFormat}
  provides a basic particle plotter which dumps particles, their cut-off radius
  plus their particle-vertex associativity. If you have to dump further
  properties, you have to write your own plotter. The generated plotter however
  can serve as a blueprint.
  \item The action set \texttt{toolbox.particles.ParticleTreeAnalysis} provides
  some fundamental tree analysis such as a multiscale counting how many
  particles reside within each cell. You'll need some additional data structures
  (per cell) to use this action set.
  \item The action set \texttt{toolbox.particles.ParticleAMR} uses the tree
  analysis to guide dynamic grid adaptivity. You can specify how many particles
  are held per cell at least, e.g., while the cut-off radius feeds naturally
  into the AMR.
\end{enumerate}



\section*{Links and further reading}

\begin{itemize}
  \item The ``official'' pidt paper is
{\tiny \begin{verbatim}
@article{WEINZIERL201642,
title = "Two particle-in-grid realisations on spacetrees",
journal = "Parallel Computing",
volume = "52",
pages = "42 - 64",
year = "2016",
issn = "0167-8191",
doi = "https://doi.org/10.1016/j.parco.2015.12.007",
url = "http://www.sciencedirect.com/science/article/pii/S0167819115001635",
author = "T. Weinzierl and B. Verleye and P. Henri and D. Roose",
keywords = "Particle-in-cell, Spacetree, Particle sorting, AMR, Lagrangian–Eulerian methods, Communication-avoiding"
}  \end{verbatim}}
  If you use the software, it would be great if you could cite this one.
\end{itemize}


